```latex
\documentclass[12pt]{article}
\usepackage[utf8]{inputenc}
\usepackage{amsmath, amssymb, graphicx, hyperref}
\usepackage{geometry}
\geometry{a4paper, margin=1in}
\title{NeuroSHOA: A Spintronic Framework for Non-Invasive Readout of Emergent Neural Dynamics}
\author{Konstantin Fedotov}
\date{\today}

\begin{document}

\maketitle

\begin{abstract}
We present NeuroSHOA, a theoretical and engineering framework that extends the SHOA formalism to biological neural systems. We prove that the Hamiltonian of neuronal microtubules is mathematically equivalent to the SHOA-Neuro Hamiltonian, establishing a direct bridge between quantum self-healing and neural dynamics. A novel readout operator $\hat{J}_{\text{readout}}(t) = \frac{d}{dt} \sum_k w_k C_k(t)$ is introduced, linking neural activity to SHOA-resilience. We propose a roadmap for non-invasive monitoring using SERF magnetometers and NV-center sensors, opening the path to devices that assess the mind's capacity for self-restoration.
\end{abstract}

\section{Introduction}
The SHOA project has evolved from quantum self-healing (Classic) through learning systems (Neuro) to adaptive materials (Adaptive). NeuroSHOA represents a logical expansion: the application of SHOA principles to biological neural networks. We argue that neuronal microtubules are a natural physical realization of a SHOA-Neuro ensemble.

\section{Model}
\subsection{Microtubule Hamiltonian}
The collective dynamics of tubulin dipoles in microtubules is described by:
\begin{equation}
\mathcal{H}_{\text{MT}} = \sum_i \frac{\delta_i}{2} \sigma_z^{(i)} + \sum_{i<j} J'_{ij} (\sigma_+^{(i)} \sigma_-^{(j)} + h.c.),
\end{equation}
where $\delta_i$ is the dipole energy gap and $J'_{ij}$ is the dipole-dipole coupling. This is mathematically identical to the SHOA-Neuro Hamiltonian.

\subsection{Readout Operator}
We introduce a new operator to non-invasively assess neural resilience:
\begin{equation}
\hat{J}_{\text{readout}}(t) = \frac{d}{dt} \sum_k w_k C_k(t),
\end{equation}
where $w_k$ are weights proportional to the SHOA-resilience of the $k$-th neural domain and $C_k(t)$ are measured cross-correlation functions. This operator not only measures activity but also evaluates the system's capacity to return to an optimal state.

\section{Experimental Roadmap}
We propose a three-stage validation:
\begin{enumerate}
    \item Laboratory (1-2 years): Magnetic field mapping of neuron slices using NV centers in diamond.
    \item Clinical (2-3 years): Development of a NeuroSHOA helmet with integrated SERF magnetometers.
    \item Theoretical: We predict that resting-state EEG patterns exhibit a ``SHOA-resilience'' signature -- the ability to restore correlations after an artificial perturbation.
\end{enumerate}

\section{Conclusion}
NeuroSHOA proves the universality of the SHOA formalism, opening a new frontier in non-invasive neural monitoring and restoration.

\end{document}
